% Clase 20 --- La analogía RC / RL (§2.2).
% No es una curiosidad estética: es la razón por la que no hay que resolver nada
% nuevo. La figura pone los dos circuitos y el diccionario entre ambos.
\begin{center}
\begin{tikzpicture}[scale=1]

  \begin{scope}
    \draw (0,0) to[battery1, l=$\varepsilon$] (0,2.2)
          to[R, l=$R$] (2.6,2.2);
    \draw (2.6,2.2) to[C, l_=$C$] (2.6,0) -- (0,0);
    \node[figlbl, anchor=north, align=center] at (1.3,-0.5)
      {(a) $RC$: $\dfrac{Q}{C} + R\dfrac{dQ}{dt} = \varepsilon$};
  \end{scope}

  \begin{scope}[xshift=6.0cm]
    \draw (0,0) to[battery1, l=$\varepsilon$] (0,2.2)
          to[R, l=$R$] (2.6,2.2);
    \draw (2.6,2.2) to[L, l_=$L$] (2.6,0) -- (0,0);
    \node[figlbl, anchor=north, align=center] at (1.3,-0.5)
      {(b) $RL$: $R\,I + L\dfrac{dI}{dt} = \varepsilon$};
  \end{scope}

  \node[figlbl, anchor=north, align=center] at (4.3,-1.7)
    {$Q \leftrightarrow I$,\qquad
     $\dfrac{1}{C} \leftrightarrow R$,\qquad
     $R \leftrightarrow L$,\qquad
     $\tau = RC \leftrightarrow \tau = \dfrac{L}{R}$};

\end{tikzpicture}\\[0.5em]
{\small\itshape Las dos ecuaciones son la misma con otros nombres, y eso vale
más que una analogía decorativa: significa que toda la matemática ya está hecha
y que basta traducir. El diccionario asigna la carga a la corriente, el inverso
de la capacitancia a la resistencia y la resistencia a la inductancia ---sí, la
$R$ cambia de papel al pasar de un circuito al otro, que es la parte
contraintuitiva del asunto---. De ahí que las soluciones tengan exactamente la
misma forma, exponenciales con una constante de tiempo, y que baste con conocer
una para escribir la otra. La lectura física, en cambio, es distinta: el
condensador almacena energía en el campo \emph{eléctrico} entre sus placas y la
bobina en el campo \emph{magnético} de su interior, y mientras el transitorio
del $RC$ se acelera cuando la resistencia baja, el del $RL$ hace lo contrario.
Un tiempo característico chico significa siempre lo mismo: el circuito llega
rápido a su régimen estacionario.}
\end{center}
